%-------------------------------------------------------------------------------
% biorag — Development History and Design Rationale
% Build:  pdflatex development-history.tex   (twice, for the ToC)
%-------------------------------------------------------------------------------
\documentclass[a4paper,11pt]{article}

\usepackage[utf8]{inputenc}
\usepackage[T1]{fontenc}
\usepackage[english]{babel}
\usepackage{charter}
\usepackage[margin=2.4cm]{geometry}
\usepackage{microtype}
\usepackage{booktabs}
\usepackage{longtable}
\usepackage{array}
\usepackage{enumitem}
\usepackage{xcolor}
\usepackage{titlesec}
\usepackage{fancyhdr}
\usepackage{tikz}
\usepackage{caption}
\usepackage[hidelinks]{hyperref}

\definecolor{accent}{HTML}{1F4E79}
\definecolor{muted}{HTML}{5A5A5A}
\definecolor{rule}{HTML}{C8C8C8}
\definecolor{boxbg}{HTML}{F4F6F8}

\hypersetup{colorlinks=true, linkcolor=accent, urlcolor=accent}

\titleformat{\section}{\large\bfseries\color{accent}}{\thesection}{0.7em}{}
\titleformat{\subsection}{\normalsize\bfseries}{\thesubsection}{0.6em}{}
\titlespacing*{\section}{0pt}{1.6em}{0.7em}

\pagestyle{fancy}
\fancyhf{}
\fancyhead[L]{\footnotesize\color{muted}biorag — Development History}
\fancyhead[R]{\footnotesize\color{muted}\thepage}
\renewcommand{\headrulewidth}{0.4pt}
\renewcommand{\headrule}{\hbox to\headwidth{\color{rule}\leaders\hrule height \headrulewidth\hfill}}

\setlist[itemize]{leftmargin=1.3em, itemsep=2pt, topsep=3pt}
\newcommand{\code}[1]{\texttt{\small #1}}

% A "decision" callout: what was chosen, and what was rejected.
\newenvironment{decision}[1]
  {\par\vspace{0.5em}\noindent
   \begin{tikzpicture}
   \node[fill=boxbg, rounded corners=2pt, inner sep=9pt, text width=\dimexpr\linewidth-20pt\relax,
         align=left] \bgroup \textbf{#1}\par\vspace{3pt}\small}
  {\egroup;\end{tikzpicture}\par\vspace{0.4em}}

%-------------------------------------------------------------------------------
\begin{document}

\begin{center}
  {\LARGE\bfseries biorag}\\[0.35em]
  {\large Development History and Design Rationale}\\[1.1em]
  {\color{muted}Retrieval-augmented QA over PubMed literature\\
   on allosteric kinase inhibitors}\\[1.4em]
  \begin{tabular}{@{}r@{\hspace{1em}}l@{}}
    \textbf{Author}     & Florian Ebner \\
    \textbf{Repository} & \url{https://github.com/FlorianEbner96/biorag} \\
    \textbf{Licence}    & MIT \\
    \textbf{Compiled}   & \today \\
  \end{tabular}
\end{center}

\vspace{1.2em}
\hrule height 0.4pt
\vspace{1.2em}

\noindent
This document records \emph{how} the system was built and, more importantly,
\emph{why} each layer was designed the way it was. It is deliberately not a
feature list: a feature list says what exists, and the interesting content of an
engineering project is the set of alternatives that were considered and
rejected. Where a decision went against the obvious default, the reasoning is
given in full, including the cases where the measured result contradicted the
expectation.

\tableofcontents

%-------------------------------------------------------------------------------
\section{The problem and the core promise}

A language model asked about allosteric kinase inhibitors will answer from its
parametric memory. In a scientific domain this is not merely imprecise; it is
unverifiable, which is worse. The reader cannot tell a correct recollection from
a fluent fabrication, because both look identical.

\begin{decision}{Core promise}
Every generated claim traces to a specific abstract via its PMID, and
\emph{every citation is checked against what was actually retrieved}. Provenance
is carried through every layer rather than reconstructed at the end, because
provenance rebuilt after the fact is provenance that is wrong sometimes.
\end{decision}

\noindent
The scope was narrowed deliberately to one target class --- allosteric kinase
inhibitors --- rather than oncology or pharmacology at large. A corpus of a few
hundred abstracts is large enough to be a genuine retrieval problem and small
enough to finish, evaluate and understand alone.

%-------------------------------------------------------------------------------
\section{Timeline}

The system was built in four commits across three days, then extended with an
HTTP layer. Each stage is listed with what it added and what question it
answered.

\begin{longtable}{@{}p{2.05cm}p{2.6cm}p{9.6cm}@{}}
\toprule
\textbf{Date} & \textbf{Commit} & \textbf{What it added} \\
\midrule
\endhead
2026-07-19 & \code{67c34f0} &
  \textbf{Foundation: data through retrieval.}
  Config, canonical schema, PubMed fetcher (NCBI E-utilities), corpus cleaner,
  chunker, embedder, vector store, and three retrievers (dense, BM25, hybrid).
  Synthetic corpus generator added at the same time, so the cleaner could be
  asserted against known ground truth from the first commit rather than
  retrofitted with tests later. \\
\addlinespace
2026-07-19 & \code{664c654} &
  \textbf{Generation and evaluation.}
  Prompt templates, the \code{Answerer} with citation verification, the gold-set
  loader and the retrieval metrics. MIT licence added. Dependencies corrected in
  the same commit: \code{numpy} and \code{requests} were used everywhere but never
  declared (a fresh install relied on them arriving transitively via
  \code{torch}), while \code{anthropic}, \code{fastapi}, \code{uvicorn} and
  \code{pydantic} were declared but unused and became optional extras.
  \code{regex} was listed and never used at all. \\
\addlinespace
2026-07-21 & \code{9cfd6eb} &
  \textbf{Query CLI and batching fix.}
  \code{examples/ask.py}: one-shot or interactive querying against the real
  index, with optional JSON output. Degrades to retrieval-only without an API
  key. \\
\addlinespace
2026-07-21 & \code{b39e488} &
  \textbf{Gold standard and evaluation runner.}
  18 hand-written questions with known relevant PMIDs, and
  \code{examples/evaluate.py} to score the three retrievers against them. \\
\addlinespace
2026-09-02 & \emph{(this work)} &
  \textbf{HTTP API.}
  \code{/health}, \code{/ask} and \code{/search} on FastAPI, with Pydantic
  request and response models, dependency injection for testability, and 17
  additional offline tests. \\
\bottomrule
\end{longtable}

\noindent
\textbf{Current size.} 2\,571 lines of source across 14 modules, 912 lines of
tests, 119 tests passing --- none requiring network access, an API key, or a
model download.

%-------------------------------------------------------------------------------
\section{Architecture}

Data flows in one direction, and every stage preserves the PMID.

\begin{center}
\small
\begin{tabular}{@{}llp{7.6cm}@{}}
\toprule
\textbf{Layer} & \textbf{Module} & \textbf{Responsibility} \\
\midrule
Data --- fetch     & \code{data/fetch\_pubmed.py}   & ESearch $\rightarrow$ EFetch, rate-limited, cached as raw JSON \\
Data --- clean     & \code{data/clean\_corpus.py}   & Dedup, boilerplate stripping, retraction and quality filters \\
Ingestion --- chunk & \code{ingestion/chunker.py}   & Section-aware splitting, whole abstracts by default \\
Ingestion --- embed & \code{ingestion/embedder.py}  & Local sentence-transformers, normalised vectors \\
Ingestion --- store & \code{ingestion/store.py}     & Exact NumPy search; optional Chroma backend \\
Retrieval          & \code{retrieval/retriever.py}  & Dense, BM25, and rank-fused hybrid \\
Generation         & \code{generation/}             & Prompting, answering, citation verification \\
Evaluation         & \code{evaluation/}             & Gold set, retrieval and refusal metrics \\
API                & \code{api/}                    & FastAPI service over the same components \\
\bottomrule
\end{tabular}
\end{center}

%-------------------------------------------------------------------------------
\section{Design decisions}

\subsection{Abstracts are not chunked by default}

The reflexive default in most RAG tutorials is a 200-token sliding window. That
convention is inherited from chunking long PDFs and is actively harmful here.

\begin{decision}{Decision}
Keep whole abstracts intact below 2\,000 characters (roughly 95\% of them).
Above that, split only on the author's own section boundaries.
\end{decision}

\noindent
A PubMed abstract is around 250 words and is already a coherent retrieval unit,
written by an expert to summarise one paper. Splitting it separates a finding
from the condition that qualifies it, and embeds fragments no model was trained
on.

\subsection{Unicode normalised, never ASCII-folded}

\begin{decision}{Decision}
NFKC normalisation only. The idiom
\code{normalize('NFKD').encode('ascii','ignore')} is explicitly rejected.
\end{decision}

\noindent
PI3K$\alpha$ and PI3K$\beta$ are different proteins with different inhibitors.
The standard cleaning idiom collapses both to ``PI3K'' and destroys the exact
distinction a question may turn on.

\subsection{Exact search, not a vector database}

\begin{decision}{Decision}
Brute-force NumPy dot product as the default store. A Chroma backend exists
behind the same interface for persistence.
\end{decision}

\noindent
At this scale a $3000 \times 384$ float32 matrix is 4.6\,MB and one
matrix--vector product takes single-digit milliseconds. Approximate-nearest-neighbour
indexes trade recall for speed at $10^5$--$10^8$ vectors; at $10^3$ they add a
dependency and lose recall in exchange for imperceptible gain.

\subsection{Hybrid retrieval fuses ranks, not scores}

\begin{decision}{Decision}
Reciprocal Rank Fusion over the two ranked lists.
\end{decision}

\noindent
Cosine similarities and BM25 scores are not commensurable, and BM25's scale is
corpus-dependent. Fusing ranks needs no score normalisation and no tuning
constant that silently depends on the corpus.

\subsection{Fuzzy title deduplication tuned tight, at 0.97}

\begin{decision}{Decision}
0.97, not the more intuitive 0.90--0.92.
\end{decision}

\noindent
A group publishing serially on related compounds produces titles sharing more
than 90\% of their characters while describing entirely different work. Testing
against a synthetic corpus with known ground truth showed that a 0.92 threshold
silently discarded roughly 15\% of genuinely distinct papers. On the real
corpus, 0.97 removed 4 near-duplicates out of 413 fetched.

\subsection{Citations are verified, not trusted}

\begin{decision}{Decision}
The prompt requests inline \code{[PMID]} markers; the answerer then checks every
cited PMID against what was actually in the context. \code{Answer.is\_grounded}
becomes false when a citation was fabricated.
\end{decision}

\noindent
This is the single most damaging failure mode the system has, precisely because
the output looks entirely credible. A model citing a plausible PMID that was
never retrieved produces text that a reader cannot distinguish from a correct
answer.

\subsection{Empty retrieval refuses deterministically}

If nothing was retrieved, the system returns a refusal without calling the model
at all. Hoping a model refuses on empty context is a coin flip; refusing in code
is free and certain.

\subsection{Everything is testable offline}

\begin{decision}{Decision}
\code{data/synthetic.py} generates PubMed-shaped records with deliberately
planted pathologies --- retractions, boilerplate, duplicate titles, one-line
abstracts. A hashing \code{FakeEmbedder} and a scripted \code{FakeLLM} stand in
for the model and the API.
\end{decision}

\noindent
Tests that need the network get skipped, and skipped tests protect nothing. All
119 tests run with no network, no API key and no model download.

%-------------------------------------------------------------------------------
\section{Evaluation}

\subsection{The gold set}

18 hand-written questions against the real corpus, written \emph{before} tuning
anything --- questions written after seeing what the system does well measure
the author's selection bias rather than the system.

\begin{center}
\small
\begin{tabular}{@{}lr@{}}
\toprule
\textbf{Question type} & \textbf{Count} \\
\midrule
Identifier (specific compound codes) & 5 \\
Factual                              & 4 \\
Conceptual                           & 4 \\
\textbf{Unanswerable} (expects refusal) & \textbf{4} \\
Multi-hop                            & 1 \\
\midrule
Total & 18 \\
\bottomrule
\end{tabular}
\end{center}

\begin{decision}{Why 22\% of the gold set is deliberately unanswerable}
This is the category most evaluation sets omit, and the only one that detects a
system which answers everything confidently. A RAG system scoring perfectly on
factual questions while never refusing is useless in a scientific setting, and
\code{refusal\_recall} cannot be computed at all without such questions.
\end{decision}

\subsection{Retrieval results}

Corpus: 409 PubMed abstracts, 494 chunks. Scored over the 14 answerable
questions at $k=5$.

\begin{center}
\small
\begin{tabular}{@{}lrrrr@{}}
\toprule
\textbf{Retriever} & \textbf{recall@5} & \textbf{precision@5} & \textbf{MRR} & \textbf{hit rate} \\
\midrule
Dense (semantic, MiniLM) & 0.274 & 0.100 & 0.339 & 0.500 \\
BM25 (keyword baseline)  & 0.405 & 0.143 & \textbf{0.393} & 0.571 \\
\textbf{Hybrid (RRF)}    & \textbf{0.429} & \textbf{0.157} & 0.392 & \textbf{0.643} \\
\bottomrule
\end{tabular}
\end{center}

\subsection{The result that contradicted the expectation}

The textbook expectation is that semantic retrieval beats keyword search. Here
it did not: dense retrieval was the \emph{weakest} of the three. Two honest
limitations explain it, and neither was known in advance.

\begin{enumerate}[leftmargin=1.6em]
  \item \textbf{The gold set leans on exact identifiers.} Nine of the 14
        answerable questions hinge on a specific compound code --- EAI045,
        ABL001, MK-2206, IACS-13909, SHP099. This is precisely what BM25 is
        built for and what a general embedding model handles poorly, since a
        compound code is a rare token it has barely seen.
  \item \textbf{The embedding model is general-purpose, not biomedical.}
        \code{all-MiniLM-L6-v2} was never trained on this vocabulary. A domain
        model such as \code{S-PubMedBert-MS-MARCO} should lift the dense numbers
        materially --- and the evaluation harness exists precisely to
        \emph{measure} that swap rather than assume it.
\end{enumerate}

\noindent
The conclusion that does hold is the one that matters: \textbf{fusing the two
beats either alone} on recall and hit rate. Neither retriever is redundant. BM25
nails the identifiers; dense contributes on the conceptual questions whose
wording shares no tokens with the abstract that answers them.

\begin{decision}{A note on BM25 and negative IDF}
On a small synthetic corpus the domain terms (``allosteric'', ``kinase'',
``inhibitor'') appear in more than half of all documents and take negative IDF,
so BM25 scores come out negative --- correctly reporting that the domain's
defining terms carry no discriminative power. On the real 409-abstract corpus
they do not, which is itself an argument for evaluating on real rather than toy
data.
\end{decision}

\subsection{Retrieval and generation are scored separately}

In that order, and on purpose. If recall@5 is 0.4, no amount of prompt
engineering will fix the answers, because the relevant abstract was never in the
context. Debugging generation when the fault is upstream in retrieval is the
most common time sink in these projects.

%-------------------------------------------------------------------------------
\section{The HTTP layer}

The API exposes the same components over HTTP without changing their contracts.

\begin{center}
\small
\begin{tabular}{@{}llp{8.2cm}@{}}
\toprule
\textbf{Endpoint} & \textbf{Method} & \textbf{Purpose} \\
\midrule
\code{/health} & GET  & Readiness, not just liveness --- reports indexed chunk count and whether generation is configured \\
\code{/ask}    & POST & Cited answer when a key is set, otherwise the retrieved abstracts \\
\code{/search} & GET  & Retrieval only, never generation --- the guaranteed-free path \\
\bottomrule
\end{tabular}
\end{center}

\begin{decision}{Load once, in a lifespan handler}
At import time the module would be unimportable without a built index, which
breaks the test suite. Per request it would pay an $\approx$80\,MB model load on
every call. FastAPI's \code{lifespan} does it exactly once and is the documented
replacement for the deprecated \code{@app.on\_event("startup")}.
\end{decision}

\begin{decision}{Dependencies are injected, never constructed in a handler}
The rest of the codebase already works this way --- \code{Answerer} takes an
injected LLM client, \code{RAGPipeline} an injected retriever --- and that is
exactly why the tests need no key, no network and no model. Handlers depend on
\code{get\_deps()}, which tests override with fake-backed instances, so the API
layer does not become the one place that breaks the property.
\end{decision}

\begin{decision}{A missing index returns 503, not 500}
``I have no index'' is not a server error, it is \emph{not ready to serve}. An
orchestrator keeps a 503 container out of rotation until its volume is mounted;
a 500 looks like a crash loop. Correspondingly, \code{/health} reports the chunk
count, because the failure that actually happens in a container is not a dead
process but a live one whose index volume was never mounted.
\end{decision}

\noindent
\textbf{The verification result is part of the response.} \code{is\_grounded}
and \code{hallucinated\_citations} are first-class response fields rather than
server-side log lines. An API returning the answer text while dropping the
verification would hand back exactly the unverifiable output this project exists
to avoid.

%-------------------------------------------------------------------------------
\section{Known limitations and next steps}

Stated plainly, because a portfolio project that claims no limitations is not
credible.

\begin{itemize}
  \item \textbf{The embedding model is not biomedical.} This is the single
        clearest improvement available, and the harness exists to measure it.
  \item \textbf{Absolute retrieval numbers are modest.} recall@5 of 0.43 on an
        identifier-heavy gold set is a starting point, not a finished result.
  \item \textbf{The gold set is small} at 18 questions, and written by one
        person. It detects gross regressions, not subtle ones.
  \item \textbf{Groundedness is checked at citation level, not claim level.} A
        model can cite a genuinely retrieved PMID for a claim that abstract does
        not make. Claim-level verification is the natural next layer.
  \item \textbf{No CI yet.} The tests run offline in under two seconds, so
        wiring them into GitHub Actions is the obvious next step and would make
        the test suite a gate rather than a habit.
  \item \textbf{Docker} is the remaining packaging step, so the whole system
        runs with one command.
\end{itemize}

\vspace{1em}
\hrule height 0.4pt
\vspace{0.8em}
\noindent{\footnotesize\color{muted}
Source and full documentation:
\url{https://github.com/FlorianEbner96/biorag}}

\end{document}
